% New eighth-hour result; root audit and independent matrix replay retained.
\subsection{An exact obstruction to identifying the two finite CM outputs}

Let $\mathcal M=\Gamma_{1,[9]}$ be the ordinary oriented
nine-point torus mapping class group, with the marked
three-division translation subgroup $H\simeq C_3^2$.
Let $\mathcal M_0$ be the subgroup fixing the distinguished
zero label.  Retain the actual surface representation
and the ordinary forgetting map
\[
 \rho:\mathcal M\longrightarrow
       G\leq\operatorname{Out}(Q_f\times Q_L),
 \qquad
 j_0:\mathcal M_0\longrightarrow\Gamma_{1,1}
                           \longrightarrow J_c
                                                               \tag{FCD.1}
\]
from A.4.25 and A.4.22.  The last map is the actual
ordinary action on $\mathcal P_c$, including its full
finite output of order $1152\cdot3^4$.
Induce $j_0$ on the nine actual cosets of $\mathcal M_0$
and retain the full label permutation.  Write
\[
 G^\dagger=operatorname{im}\left(
 \mathcal M\longrightarrow
 G\times\bigl(J_c^9\rtimes\mathfrak S_9\bigr)
                    \times\mathfrak S_9\right),
 \qquad H^\dagger=\operatorname{im}(H).
                                                               \tag{FCD.2}
\]
The convention for induction is the actual coset action
with its source-word weights; at the identity coset it
recovers $j_0$ on $\mathcal M_0$.  Thus it defines an
actual homomorphism $j_0^\dagger$ on the image of
$\mathcal M_0$ in $G^\dagger$.

\begin{result}{A.4.42}{Proved}{The joined finite normalizer has an order-three discrepancy between descent and forgetting}
There is an element
\[
 n\in C_{G^\dagger}(H^\dagger)
                     \cap\operatorname{im}(\mathcal M_0)
                                                               \tag{FCD.3}
\]
of exact order three such that
\[
       j_0^\dagger(n)=1,\qquad
       d\bigl(\operatorname{pr}_G(n)\bigr)\ne1.
                                                               \tag{FCD.4}
\]
The latter value also has exact order three.  Here $d$
is precisely the homomorphism (TF.4), with its specified
source $N_G(H)$ and target $J_c$.

The same assertion holds after adjoining any finite
family of actual finite source coordinates whose joint
image on the pure quaternion-kernel domain of A.4.23 is
solvable.  In particular every finite family of additional
solvable-valued coordinates leaves such a discrepancy.
This statement gives the exact failure of the proposed
equality of these two maps on this finite normalizer.
Their equality on the image of the actual ordinary
normalizer remains the comparison already proved in
A.4.25.
\end{result}

\begin{proof}
\par\medskip\noindent\textbf{An explicit order-three centralizer element.}\quad
Use the actual eighteen-dimensional multiplicity space
$(M,F)$ and its deck matrices $h_1,h_2$ from A.4.15.
Let $R_v$ be the first of the two explicitly computed
Hermitian reflections of A.4.20, with augmentation norm
one.  Its negative eigenspace is the nondegenerate
nine-dimensional free $\mathbb F_3[H]$-line $Av$.
The actual orthogonal projector onto this line is
\[
          P=\frac{I-R_v}{2},\qquad
          k=h_1P+h_1^{-1}(I-P).
                                                               \tag{FCD.5}
\]
The two eigenspaces are orthogonal and $H$-stable.
The matrices $P,h_1,h_2$ commute, $P^2=P$, and
$P^tF=FP$.  Hence $k$ preserves $F$, centralizes $H$,
and has order dividing three.  Its displayed nontrivial
quotient below proves that its order is exactly three.
Both determinant and spinor norm take values in groups
of order two.  Their values on this odd-order isometry
are therefore one, proving
\[
                    k\in C_{\Omega_{18}^-(3)}(H).
                                                               \tag{FCD.6}
\]

Let $\varphi:M\to M_1=M/\mathfrak m_H^2M$ be the
actual six-dimensional quotient map of A.4.21.  In its
retained free $A_1$ basis, ordered by degree zero,
the $X=h_1-1$ multiples and the $Y=h_2-1$ multiples,
the induced matrix is exactly
\[
 \overline k=
 \begin{pmatrix}
  I_2&0&0\\
  D&I_2&0\\
  0&0&I_2
 \end{pmatrix},
 \qquad D=\begin{pmatrix}0&1\\1&0\end{pmatrix}.
                                                               \tag{FCD.7}
\]
These are the same fixed quotient coordinates used by
the complete ordinary-image calculation in A.4.21.
The inner kernel in the compatible-pair reconstruction
of A.4.22 consists of the eighteen module values
$\{\pm h:h\in H\}$.
An element of this set with constant term $I_2$ has
first-degree coefficients $aI_2,bI_2$ for some
$a,b\in\mathbb F_3$.  The coefficient $D$ in (FCD.7)
is not scalar.  Thus $\overline k$ lies outside this
entire inner kernel.  The compatible pair with this
module action and identity base action gives a
nontrivial outer automorphism of $\mathcal P_c$.

\par\medskip\noindent\textbf{Realization with every ordinary forgetting value fixed.}\quad
Let $\mathcal G$ be the group of actual based pure
quaternion-kernel geometric automorphisms in A.4.23.
Its multiplicity map
$\rho_M:\mathcal G\twoheadrightarrow
\mathrm{SO}_{18}^-(3)$ is onto.  Retain its actual
lower map $\beta$ from (PC.2), whose image is solvable
of derived length at most five.

On $\mathcal G$, the induced $j_0$ representation has
trivial permutation part and takes values in $J_c^9$.
This is a solvable group.  More explicitly, the
modulo-two kernel of the ordinary image is a quotient
of the group in (NJ.7).  That group has an elementary
abelian $3^6$ kernel over a subgroup of $C_4\times P$,
where $P=C_2^3\times\mathrm{SL}_2(\mathbb F_3)$
in the Chinese-remainder coordinates of (NJ.5).
The displayed series
$\mathrm{SL}_2(\mathbb F_3),Q_8,C_2,1$
therefore gives derived length at most four for that
modulo-two kernel.  Restoring the quotient
$\mathrm{SL}_2(\mathbb F_2)\simeq\mathfrak S_3$
gives derived length at most six for $J_c$.
Its ninth power has the same bound.

Form the actual finite joint image
\[
 E=(\rho_M,\beta,\operatorname{Ind}j_0)(\mathcal G)
 \leq\mathrm{SO}_{18}^-(3)\times\mathcal B\times J_c^9.
                                                               \tag{FCD.8}
\]
Projection of a derived subgroup is the derived subgroup
of the projected image: lift the two entries of each
commutator and multiply their lifts.  A.4.23 proves
$\mathrm{SO}_{18}^-(3)^{(r)}=\Omega_{18}^-(3)$ for
every $r\geq1$.  Thus the sixth derived subgroup of
(FCD.8) projects onto $\Omega_{18}^-(3)$, while both
remaining projections are identity.  It follows that
\[
                  (k,1,1)\in E^{(6)}\subset E.
                                                               \tag{FCD.9}
\]
This is membership in the actual source image under
the displayed homomorphism.  Consequently a finite
word in actual pure quaternion-kernel mapping classes
realizes $k$, has all lower coordinates (PC.2) equal
to identity, and has every one of its nine ordinary
forgetting values equal to identity.

For that word the joint quaternion and compact-mod-four
coordinates give identity base action on $G_2$.
The only remaining ambiguity in its faithful factor
is the inner module ambiguity of A.4.19.  Its outer
faithful coordinate is therefore exactly the compatible
pair $(k,1)$.  The three lower identities in (PC.2)
give identity outer action on $Q_L$ by (JP.7), as in
the proof of A.4.25.  Hence its image in (FCD.2) is
the fully specified element with compatible faithful
coordinate $(k,1)$ and all other coordinates identity.
Call this element $n$.

It is pure, so fixes the zero label.  It centralizes
the faithful deck action because $k$ commutes with
both actual $H$ matrices; compatible pairs determine
outer classes by A.4.19.  Its $Q_L$ coordinate and
its whole induced-forgetting coordinate are identity,
so those commute with $H^\dagger$ as well.  This proves
(FCD.3).  Its identity-coset forgetting value is one.
Formula (TF.15a) sends its faithful component to the
compatible pair $(\overline k,1)$, whose outer class
is nontrivial by (FCD.7).  Every coordinate of $n$
has order dividing three and this last image is
nontrivial.  Thus both asserted orders and (FCD.4)
follow.

For any additionally specified finite family in the
statement, join its actual solvable image to the last
two coordinates of (FCD.8).  Choose the derived depth
at least its finite derived length and at least six.
The same projection identity and perfectness give
$(k,1,\ldots,1)$ in that actual joint image.  Its
additional coordinates are identity and consequently
centralize the added deck images.  This proves the
assertion for the full stated family, with the source
membership and every projection retained.

The executed program
\path{code/check_finite_cm_forgetting_discrepancy.py}
checks the actual reflection and rank-nine projector,
all orthogonality and deck-commutation equations,
the exact order three, the quotient intertwining
$\varphi k=\overline k\varphi$, and all eighteen
inner-kernel comparisons.  It prints (FCD.7) in the
retained basis and archives every matrix and input
hash.  An independent fresh replay gives the same
coefficients.  The actual source membership is the
proved joint-image identity (FCD.9); it is not inferred
from a matrix order or an ambient dimension count.
\end{proof}
